%% Carta de presentacion para Computers & Industrial Engineering. Lleva
%% identidad: no es el manuscrito anonimo. La firma el autor.
\documentclass[11pt]{article}
\usepackage[margin=1.9cm]{geometry}
\usepackage{lmodern}
\usepackage[T1]{fontenc}
\usepackage[utf8]{inputenc}
\usepackage{parskip}
\setlength{\parskip}{0.4em}
\pagestyle{empty}

\begin{document}

\noindent
Hern\'an D\'iaz Rodr\'iguez\\
Department of Computing, University of Oviedo\\
Gij\'on, Spain\\
diazhernan@uniovi.es

\vspace{1em}
\noindent 5 October 2026

\vspace{1em}
\noindent To the Editors,\\
\emph{Computers \& Industrial Engineering}

\vspace{1em}

I submit for your consideration the manuscript \emph{Genetic
Programming Hyper-Heuristics for the Job Shop Scheduling Problem with
Interval Durations: Robustness-Aware Interpretable Rules that Generalize
Across Instance Sizes}, as a research paper.

In many manufacturing settings the duration of an operation is not
known in advance, and a planner can only give a lower and an upper
bound. The interval job shop scheduling problem models this situation,
and its best solutions so far come from metaheuristics that search anew
for every instance. The manuscript brings automated rule design to this
problem. A genetic programming hyper-heuristic evolves dispatching rules
that read both the worst-case value and the width of each uncertain
quantity. The rules build a schedule in milliseconds and are short
enough to be read and checked by a planner.

The methodological contribution has three parts. First, rules evolved
on four instances improve on every constructive baseline on all seventy
instances of the benchmark, and transfer without retraining to all size
classes and to a second benchmark. Second, ablations show where interval
information pays off: not in the expected makespan, but in the
predictability of the schedule. An objective that penalizes the width of
the predicted makespan yields schedules whose execution departs less
from the prediction, with a single weight that sets the trade-off.
Third, a comparison at matched budget with a genetic algorithm,
reimplemented on the same decoder, shows up to which computational
budget the evolved rule is the better choice, and how that budget grows
with instance size.

The work fits the journal because it answers questions that a
production planner faces: which rule to deploy when schedules must be
built or rebuilt quickly, how much expected makespan to give up for a
more reliable schedule, and when a per-instance search is worth its
cost.

I confirm that the manuscript is original, that it has not been
published elsewhere, and that it is not under consideration by another
journal. An earlier version of this study is available as a preprint
(SSRN, doi:10.2139/ssrn.7214286).

The complete source code is openly available in Zenodo
(doi:10.5281/zenodo.21716972; version 2.0, which accompanies this
manuscript: doi:10.5281/zenodo.23156759): the hyper-heuristic, the baselines and
the genetic algorithm, together with a test that re-derives the
deposited results from scratch. The deposit also holds the instances,
the evolved rules and the result files behind every table and figure. A
statement on the use of generative AI in the preparation of the
manuscript appears directly before the references, as the journal's
policy requires. In keeping with the double-anonymous review process,
the author's identity appears only in the title page and in this letter.

Thank you for considering this work.

\vspace{0.7em}
\noindent Yours sincerely,

\vspace{1.2em}
\noindent Hern\'an D\'iaz Rodr\'iguez\\
Corresponding author

\end{document}
